\section{AIME II 2007}
        \begin{enumerate}
        	\item (\textit{AIME II 2007}) A mathematical organization is producing a set of commemorative license plates. Each plate contains a sequence of five characters chosen from the four letters in AIME and the four digits in $2007$. No character may appear in a sequence more times than it appears among the four letters in AIME or the four digits in $2007$. A set of plates in which each possible sequence appears exactly once contains N license plates. Find N/10.


 

	\item (\textit{AIME II 2007}) Find the number of ordered triples $(a,b,c)$ where $a$, $b$, and $c$ are positive \href{https://artofproblemsolving.com/wiki/index.php?title=Integer}{integers}, $a$ is a \href{https://artofproblemsolving.com/wiki/index.php?title=Factor}{factor} of $b$, $a$ is a factor of $c$, and $a+b+c=100$.


 

	\item (\textit{AIME II 2007}) \href{https://artofproblemsolving.com/wiki/index.php?title=Square}{Square} $ABCD$ has side length $13$, and \href{https://artofproblemsolving.com/wiki/index.php?title=Point}{points} $E$ and $F$ are exterior to the square such that $BE=DF=5$ and $AE=CF=12$. Find $EF^{2}$.


 
                    \begin{figure}[H]
                        \centering
                        \includegraphics[width=0.5\textwidth]{Problem/AIME_II_2007_3_Problem_Figure_1.png}
                    \end{figure}
                    

	\item (\textit{AIME II 2007}) The workers in a factory produce widgets and whoosits. For each product, production time is \href{https://artofproblemsolving.com/wiki/index.php?title=Constant}{constant} and identical for all workers, but not necessarily equal for the two products. In one hour, $100$ workers can produce $300$ widgets and $200$ whoosits. In two hours, $60$ workers can produce $240$ widgets and $300$ whoosits. In three hours, $50$ workers can produce $150$ widgets and $m$ whoosits. Find $m$.


 

	\item (\textit{AIME II 2007}) The \href{https://artofproblemsolving.com/wiki/index.php?title=Graph}{graph} of the \href{https://artofproblemsolving.com/wiki/index.php?title=Equation}{equation} $9x+223y=2007$ is drawn on graph paper with each \href{https://artofproblemsolving.com/wiki/index.php?title=Square}{square} representing one \href{https://artofproblemsolving.com/wiki/index.php?title=Unit\_square}{unit} in each direction. How many of the $1$ by $1$ graph paper squares have interiors lying entirely below the graph and entirely in the first \href{https://artofproblemsolving.com/wiki/index.php?title=Quadrant}{quadrant}?


 

	\item (\textit{AIME II 2007}) An integer is called \textit{parity-monotonic} if its decimal representation $a_{1}a_{2}a_{3}\cdots a_{k}$ satisfies $a_{i}<a_{i+1}$ if $a_{i}$ is \href{https://artofproblemsolving.com/wiki/index.php?title=Odd}{odd}, and $a_{i}>a_{i+1}$ if $a_{i}$ is \href{https://artofproblemsolving.com/wiki/index.php?title=Even}{even}. How many four-digit parity-monotonic integers are there? 


 

	\item (\textit{AIME II 2007}) Given a \href{https://artofproblemsolving.com/wiki/index.php?title=Real\_number}{real number} $x,$ let $\lfloor x \rfloor$ denote the \href{https://artofproblemsolving.com/wiki/index.php?title=Floor\_function}{greatest integer} less than or equal to $x.$ For a certain \href{https://artofproblemsolving.com/wiki/index.php?title=Integer}{integer} $k,$ there are exactly $70$ positive integers $n_{1}, n_{2}, \ldots, n_{70}$ such that $k=\lfloor\sqrt[3]{n_{1}}\rfloor = \lfloor\sqrt[3]{n_{2}}\rfloor = \cdots = \lfloor\sqrt[3]{n_{70}}\rfloor$ and $k$ divides $n_{i}$ for all $i$ such that $1 \leq i \leq 70.$


 Find the maximum value of $\frac{n_{i}}{k}$ for $1\leq i \leq 70.$


 

	\item (\textit{AIME II 2007}) A \href{https://artofproblemsolving.com/wiki/index.php?title=Rectangle}{rectangular} piece of paper measures 4 units by 5 units. Several \href{https://artofproblemsolving.com/wiki/index.php?title=Line}{lines} are drawn \href{https://artofproblemsolving.com/wiki/index.php?title=Parallel}{parallel} to the edges of the paper. A rectangle determined by the \href{https://artofproblemsolving.com/wiki/index.php?title=Intersection}{intersections} of some of these lines is called \textit{basic} if 


 
            \begin{enumerate}[label=(\roman*)]
                \item  all four sides of the rectangle are segments of drawn line segments, and
\item  no \href{https://artofproblemsolving.com/wiki/index.php?title=Segment}{segments} of drawn lines lie inside the rectangle.
            \end{enumerate}
            Given that the total length of all lines drawn is exactly 2007 units, let $N$ be the maximum possible number of basic rectangles determined. Find the \href{https://artofproblemsolving.com/wiki/index.php?title=Remainder}{remainder} when $N$ is divided by 1000.


 

	\item (\textit{AIME II 2007}) \href{https://artofproblemsolving.com/wiki/index.php?title=Rectangle}{Rectangle} $ABCD$ is given with $AB=63$ and $BC=448.$ Points $E$ and $F$ lie on $AD$ and $BC$ respectively, such that $AE=CF=84.$ The \href{https://artofproblemsolving.com/wiki/index.php?title=Inscribed\_circle}{inscribed circle} of \href{https://artofproblemsolving.com/wiki/index.php?title=Triangle}{triangle} $BEF$ is \href{https://artofproblemsolving.com/wiki/index.php?title=Tangent}{tangent} to $EF$ at point $P,$ and the inscribed circle of triangle $DEF$ is tangent to $EF$ at \href{https://artofproblemsolving.com/wiki/index.php?title=Point}{point} $Q.$ Find $PQ.$


 

	\item (\textit{AIME II 2007}) Let $S$ be a \href{https://artofproblemsolving.com/wiki/index.php?title=Set}{set} with six \href{https://artofproblemsolving.com/wiki/index.php?title=Element}{elements}. Let $\mathcal{P}$ be the set of all \href{https://artofproblemsolving.com/wiki/index.php?title=Subset}{subsets} of $S.$ Subsets $A$ and $B$ of $S$, not necessarily distinct, are chosen independently and at random from $\mathcal{P}$. The \href{https://artofproblemsolving.com/wiki/index.php?title=Probability}{probability} that $B$ is contained in one of $A$ or $S-A$ is $\frac{m}{n^{r}},$ where $m$, $n$, and $r$ are \href{https://artofproblemsolving.com/wiki/index.php?title=Positive}{positive} \href{https://artofproblemsolving.com/wiki/index.php?title=Integer}{integers}, $n$ is \href{https://artofproblemsolving.com/wiki/index.php?title=Prime}{prime}, and $m$ and $n$ are \href{https://artofproblemsolving.com/wiki/index.php?title=Relatively\_prime}{relatively prime}. Find $m+n+r.$ (The set $S-A$ is the set of all elements of $S$ which are not in $A.$)


 

	\item (\textit{AIME II 2007}) Two long \href{https://artofproblemsolving.com/wiki/index.php?title=Cylinder}{cylindrical} tubes of the same length but different \href{https://artofproblemsolving.com/wiki/index.php?title=Diameter}{diameters} lie \href{https://artofproblemsolving.com/wiki/index.php?title=Parallel}{parallel} to each other on a \href{https://artofproblemsolving.com/wiki/index.php?title=Plane}{flat surface}. The larger tube has \href{https://artofproblemsolving.com/wiki/index.php?title=Radius}{radius} $72$ and rolls along the surface toward the smaller tube, which has radius $24$. It rolls over the smaller tube and continues rolling along the flat surface until it comes to rest on the same point of its \href{https://artofproblemsolving.com/wiki/index.php?title=Circumference}{circumference} as it started, having made one complete revolution. If the smaller tube never moves, and the rolling occurs with no slipping, the larger tube ends up a \href{https://artofproblemsolving.com/wiki/index.php?title=Distance}{distance} $x$ from where it starts. The distance $x$ can be expressed in the form $a\pi+b\sqrt{c},$ where $a,$ $b,$ and $c$ are \href{https://artofproblemsolving.com/wiki/index.php?title=Integer}{integers} and $c$ is not divisible by the \href{https://artofproblemsolving.com/wiki/index.php?title=Square}{square} of any \href{https://artofproblemsolving.com/wiki/index.php?title=Prime}{prime}. Find $a+b+c.$


 

	\item (\textit{AIME II 2007}) The increasing \href{https://artofproblemsolving.com/wiki/index.php?title=Geometric\_sequence}{geometric sequence} $x_{0},x_{1},x_{2},\ldots$ consists entirely of \href{https://artofproblemsolving.com/wiki/index.php?title=Integer}{integral} powers of $3.$ Given that


 $\sum_{n=0}^{7}\log_{3}(x_{n}) = 308$ and $56 \leq \log_{3}\left ( \sum_{n=0}^{7}x_{n}\right ) \leq 57,$


 find $\log_{3}(x_{14}).$


 

	\item (\textit{AIME II 2007}) A \href{https://artofproblemsolving.com/wiki/index.php?title=Triangle}{triangular} \href{https://artofproblemsolving.com/wiki/index.php?title=Array}{array} of \href{https://artofproblemsolving.com/wiki/index.php?title=Square}{squares} has one square in the first row, two in the second, and in general, $k$ squares in the $k$th row for $1 \leq k \leq 11.$ With the exception of the bottom row, each square rests on two squares in the row immediately below (illustrated in given diagram). In each square of the eleventh row, a $0$ or a $1$ is placed. Numbers are then placed into the other squares, with the entry for each square being the sum of the entries in the two squares below it. For how many initial distributions of $0$'s and $1$'s in the bottom row is the number in the top square a \href{https://artofproblemsolving.com/wiki/index.php?title=Multiple}{multiple} of $3$?


 
                    \begin{figure}[H]
                        \centering
                        \includegraphics[width=0.5\textwidth]{Problem/AIME_II_2007_13_Problem_Figure_1.png}
                    \end{figure}
                    


 

	\item (\textit{AIME II 2007}) Let $f(x)$ be a \href{https://artofproblemsolving.com/wiki/index.php?title=Polynomial}{polynomial} with real \href{https://artofproblemsolving.com/wiki/index.php?title=Coefficient}{coefficients} such that $f(0) = 1,$ $f(2)+f(3)=125,$ and for all $x$, $f(x)f(2x^{2})=f(2x^{3}+x).$ Find $f(5).$


 

	\item (\textit{AIME II 2007}) Four \href{https://artofproblemsolving.com/wiki/index.php?title=Circle}{circles} $\omega,$ $\omega_{A},$ $\omega_{B},$ and $\omega_{C}$ with the same \href{https://artofproblemsolving.com/wiki/index.php?title=Radius}{radius} are drawn in the interior of \href{https://artofproblemsolving.com/wiki/index.php?title=Triangle}{triangle} $ABC$ such that $\omega_{A}$ is \href{https://artofproblemsolving.com/wiki/index.php?title=Tangent}{tangent} to sides $AB$ and $AC$, $\omega_{B}$ to $BC$ and $BA$, $\omega_{C}$ to $CA$ and $CB$, and $\omega$ is \href{https://artofproblemsolving.com/wiki/index.php?title=Externally\_tangent}{externally tangent} to $\omega_{A},$ $\omega_{B},$ and $\omega_{C}$. If the sides of triangle $ABC$ are $13,$ $14,$ and $15,$ the radius of $\omega$ can be represented in the form $\frac{m}{n}$, where $m$ and $n$ are \href{https://artofproblemsolving.com/wiki/index.php?title=Relatively\_prime}{relatively prime} positive integers. Find $m+n.$


 

\end{enumerate}